\documentclass[11pt]{article}
\usepackage[utf8]{inputenc}
\usepackage[T1]{fontenc}
\usepackage[margin=1in]{geometry}
\usepackage{hyperref}
\usepackage{enumitem}
\hypersetup{colorlinks=true,linkcolor=black,urlcolor=blue,citecolor=blue}
\title{Tinkercad Design and Revision}
\author{Piter Garcia}
\date{March 20, 2026}
\begin{document}
\maketitle
\section*{Packet Use}
This printable lesson packet is generated from the paired Markdown guide. The Markdown version remains the clickable, neurodivergent-friendly working copy; this PDF carries the same lesson substance in a formal handout format. Evidence claims in this packet are based on transcript-derived lesson notes and the cleaned transcript files in this workspace.

\section*{Quick Scan}
\begin{itemize}[leftmargin=*]
\item \textbf{Grade:} Grade 5
\item \textbf{Grade Hub:} Notes [../grade-5-Notes.md]
\item \textbf{Schedule Reference:} Spring2026\_TeachingSchedule\_Derek.md [../../school/Spring2026\_TeachingSchedule\_Derek.md]
\item \textbf{Workbook Sheet:} \texttt{5th} in \texttt{Teaching\_Placement-IGNITE Curriculum\_ GCSD25-26.xlsx}
\item \textbf{Lesson Focus:} Students build 3D models in Tinkercad by placing shapes, using workplanes, grouping objects, and revising designs through iterative testing.
\item \textbf{CS Alignment Strength:} Moderate CS / computational artifact design
\item \textbf{LaTeX Source:} tinkercad-design-and-revision.tex [./tinkercad-design-and-revision.tex]
\item \textbf{Bib File:} tinkercad-design-and-revision.bib [./tinkercad-design-and-revision.bib]
\item \textbf{Artifacts Folder:} artifacts/ [./artifacts/]
\end{itemize}

\section*{Lesson Identity}
\begin{itemize}[leftmargin=*]
\item \textbf{Likely Class Blocks:} Day 2 -- 10:45-11:35 -- Pickett, Day 3 -- 10:45-11:35 -- McGlashon, Day 5 -- 10:45-11:35 -- Pecora
\item \textbf{Main Lesson Family:} Tinkercad Design and Revision
\item \textbf{Primary Evidence Base:} Transcript-derived notes plus source transcripts from this teaching-placement workspace.
\item \textbf{Primary External Source Type:} Official curriculum/provider materials.
\end{itemize}

\section*{What We Are Teaching}
\begin{itemize}[leftmargin=*]
\item navigating a digital design workspace
\item breaking a design into component shapes
\item using grouping and holes to refine an artifact
\item revising a computational artifact after testing or feedback
\end{itemize}

\section*{CS Standards Alignment}
\begin{itemize}[leftmargin=*]
\item \textbf{1B-AP-11}: Decompose problems into smaller, manageable subproblems. Why it fits here: Students break a complex object into shapes and grouped steps.
\item \textbf{1B-AP-13}: Use an iterative process to plan the development of a computational artifact. Why it fits here: Students revise, regroup, and refine the design after testing how it looks or works.
\item \textbf{1B-IC-20}: Seek diverse perspectives for the purpose of improving computational artifacts. Why it fits here: Students benefit from feedback on whether the artifact is understandable and functional.
\end{itemize}

\section*{NYS / Local Curriculum Alignment}
\begin{itemize}[leftmargin=*]
\item Aligned to the local Grade 5 design sequence on the workbook sheet through digital modeling and revision.
\item Supports New York-facing technology and design outcomes by requiring students to create and iteratively improve a digital artifact.
\item This packet is less programming-heavy but still part of computing because students use structured digital design processes and iterative computational artifact creation.
\end{itemize}

\section*{Lesson Content and Concepts}
\begin{itemize}[leftmargin=*]
\item computational artifact
\item decomposition
\item grouping
\item iteration
\item spatial planning
\end{itemize}

\section*{Evidence / Proof of Instruction}
\begin{itemize}[leftmargin=*]
\item \textbf{Transcript-derived note:} 260304 Grade 5 Pickett Tinkercad [../../transcript-derived-lessons/260304-grade-5-pickett-tinkercad.md] The note explicitly lists shape placement, planes, grouping, and hole creation as taught content.
\item \textbf{Source transcript:} 260304 Grade 5 Pickett Tinkercad [../../transcripts/260304-grade-5-pickett-tinkercad.md] The source transcript documents classroom guidance around how to enter the workspace and manipulate design tools.
\item \textbf{Observed teacher moves:} Instruction is evidenced by modeling the design workspace, sequencing tool use, and revising models rather than leaving students to explore randomly.
\item \textbf{Likely student proof:} Saved Tinkercad models, screenshots of revisions, and student explanations of grouping or hole use count as proof.
\end{itemize}

\section*{Assessment Evidence}
\begin{itemize}[leftmargin=*]
\item Student can place and manipulate at least one basic shape intentionally.
\item Student can use grouping or holes to improve a design.
\item Student can explain one revision they made and why it improved the artifact.
\end{itemize}

\section*{UDL and Inclusion Supports}
\subsection*{Multiple Representation}
\begin{itemize}[leftmargin=*]
\item Project the workspace and name only the tools needed for the current task.
\item Use side-by-side before/after screenshots to show revision value.
\item Keep navigation vocabulary tied to visible actions.
\end{itemize}

\subsection*{Multiple Action / Expression}
\begin{itemize}[leftmargin=*]
\item Allow students to sketch first, then build.
\item Provide a starter file or partially built model for students who need lower entry load.
\item Accept oral explanation of design choices in addition to written reflection.
\end{itemize}

\subsection*{Multiple Engagement}
\begin{itemize}[leftmargin=*]
\item Use real object examples so the design has a concrete target.
\item Frame revision as improvement, not mistake-fixing.
\item Let students compare two design choices and defend one.
\end{itemize}

\section*{How We Are Already Making It Inclusive}
\begin{itemize}[leftmargin=*]
\item Transcript evidence shows the lesson was already focused on tool-specific instruction, not unstructured exploration.
\item The workspace lends itself to strong visual modeling, which supports many learners.
\item Iterative revision naturally creates opportunities for students to succeed through multiple attempts.
\end{itemize}

\section*{How to Improve Inclusion Next Time}
\begin{itemize}[leftmargin=*]
\item Add a printed icon guide for move, rotate, group, hole, and workplane.
\item Capture before/after screenshots in \texttt{artifacts/} to prove revision.
\item Create a two-level challenge: basic shape composition and advanced grouped design.
\item Add a quick oral checkpoint where students explain their design plan before building.
\end{itemize}

\section*{Materials and Setup}
\begin{itemize}[leftmargin=*]
\item Student devices
\item Tinkercad accounts or access
\item Projected workspace
\item Reference objects or sample images
\end{itemize}

\section*{Teacher Moves / Script / Routines}
\begin{itemize}[leftmargin=*]
\item Model one object by decomposing it into shapes before touching the software tools.
\item Pause after the first shape placement to make sure all students are in the same view.
\item Show grouping and hole tools as solutions to a design problem, not as disconnected features.
\item Close with a revision share-out.
\end{itemize}

\section*{Common Errors and Debugging}
\begin{itemize}[leftmargin=*]
\item Students may get lost in navigation and view control before understanding design concepts.
\item Grouping too early can make editing harder.
\item Students may overbuild without checking whether the design still matches the goal.
\end{itemize}

\section*{Extensions and Simplifications}
\begin{itemize}[leftmargin=*]
\item Extension: redesign the object for a different user need or constraint.
\item Simplification: provide a three-shape target with fixed measurements.
\item Extension: compare two design plans and choose the more efficient build process.
\end{itemize}

\section*{Related Local Materials}
\begin{itemize}[leftmargin=*]
\item No direct local lesson packet linked yet.
\end{itemize}

\section*{Artifacts Checklist}
\begin{itemize}[leftmargin=*]
\item Compiled lesson PDF in \texttt{artifacts/}
\item At least one screenshot, student example, or setup photo
\item One short assessment or observation record
\item Updated standards/source bibliography once the \texttt{.bib} file is revised
\item Any teacher reflection or revision notes after reteaching
\end{itemize}

\section*{Selected Official Sources}
Selected official sources that informed this packet are listed below.
ocite{*}

\bibliographystyle{plain}
\bibliography{tinkercad-design-and-revision}
\end{document}
